%! Designing Experiments — Unit 3, Lesson 3.4 — Guided Notes & Practice (Answer Key)
% MIRROR of the blank: -key for -boxes, the header, \termblank -> \termans, \blank -> \ans, and
% the fourth argument of every \pcell filled. Nothing else differs --- keep it that way.
% THE in-class document, in the MAIN IDEAS / NOTES shape (see LESSON_SHAPE.md §1), at the
% COMPACT budget: the vocabulary box, then ONE guidednotes table of FOUR rows (three instruction
% rows + Guided Practice), 14 problems, 3 pages.
%   I DO   : the teacher fills the definition lines and works the FIRST problem of each grid
%            (1, 5, 8).
%   WE DO  : the class works the rest; the last row, Guided Practice, is worked together.
%   YOU DO : nothing on this page — ../homework/ IS the individual practice.
% THE TRAP is problem 7 (block on birth month — a variable unrelated to the response).
% THE CRUX is problem 10, in the last instruction row: "each runner wears both soles, so the
% pairing already controls everything — skip the coin flip." Blocking/pairing does not replace
% random assignment (EK VAR-3.C.7, VAR-3.C.9).
%
% THE CONTEXTS (instruction, guided practice, and homework each use a different one)
%  Row 1 — gum and memory: 60 volunteers, mint gum / fruit gum / no gum, words recalled.
%  Row 2 — a sports drink and mile times, blocked on varsity vs gym-class runners (pre-drawn).
%  Row 3 — a new shoe sole (each runner wears one of each) and typists with/without music.
%  Row 4 (GP) — a zinc lozenge for colds, placebo control, blocked on age (pre-drawn).
% Inference and scope of inference belong to Lesson 3.5 — not taught here.
%
% PAGE LOCKSTEP with notes_key/: \termblank -> \termans, \blank -> \ans, 4th arg of \pcell.
% TABLE MECHANICS: inside a cell `\\` ENDS THE ROW — use \par. Long rows are split into
% sub-rows (a bare `\\` then `& ...`) so the table can break between them.
% NO SKETCH QUESTIONS: every design diagram is pre-drawn in TikZ and read, never drawn.
\documentclass[12pt]{article}
\usepackage{apstats-article}
\usepackage{apstats-key}   % pulls in -boxes; do NOT also load -boxes

\tikzset{dbox/.style={draw=navy, fill=sky, rounded corners=2pt, align=center,
  inner sep=2.5pt, font=\scriptsize}, darr/.style={->, >=stealth, navy, thick}}

\begin{document}

\pageheader{Unit 3, Lesson 3.4}{Guided Notes \& Practice --- Answer Key}
% NAMESTRIP: no \namedateperiod here — the name row lives on the cover only.

\vspace{2pt}

\boxguard[16]
\begin{vocabbox}
\small
Fill in each term \textbf{as we name it} in the notes below.
\par
\termans{Confounding variable}{a variable tied to the explanatory variable that also affects the response, so the two effects get mixed up.}
\par
\termans{Control group}{units given no treatment or a placebo, as a baseline to compare against.}
\par
\termans{Placebo}{an inactive look-alike treatment; the \emph{placebo effect} is a response to it.}
\par
\termans{Single- / double-blind}{single: one side does not know who gets which treatment; double: neither subjects nor staff know.}
\par
\end{vocabbox}

\small
\begin{guidednotes}
% ── ROW 1: the parts and the four principles ─────────────────────────────────
\mainidea[What makes it an]{Experiment} &
In an experiment the researcher \ans{assigns} treatments. The \textbf{experimental units} are
the individuals that are \ans{assigned} the treatments (\emph{subjects} when they are people).
A \textbf{factor} is an explanatory variable whose levels are set on purpose; each level (or
combination of levels) is a \textbf{treatment}. The \textbf{response variable} is measured
\ans{after} the treatments.
\par\vspace{4pt}
\textbf{Four principles of a well-designed experiment}\par\vspace{2pt}
\stepnum{1}\ \textbf{Comparison} of at least \ans{two} treatment groups (one may be a
\textbf{control group}).\par
\stepnum{2}\ \textbf{Random assignment}: \ans{chance} decides which unit gets which treatment,
so \textbf{confounding variables} tend to balance out across the groups.\par
\stepnum{3}\ \textbf{Replication}: more than \ans{one} unit in each treatment group.\par
\stepnum{4}\ \textbf{Control}: keep every other variable the \ans{same} for all groups.
\\
& A psychologist recruits 60 student volunteers. A computer randomly assigns 20 to chew mint
gum, 20 to chew fruit gum, and 20 to chew no gum while they study a list of 30 words. Then
she records how many words each student recalls.
\notesprompt{Name the parts, then the principles.}
\begin{probgrid}{|Y|Y|}
\pcell{1}{Name the experimental units and the response variable.}{1.3cm}{60 student volunteers; number of words recalled.} &
\pcell{2}{Name the factor and the treatments.}{1.3cm}{Gum: mint gum, fruit gum, no gum (3 treatments).} \\ \hline
\pcell{3}{Why should a computer, not the psychologist, decide who chews gum?}{1.5cm}{Chance spreads memory skill and study habits evenly across the three groups.} &
\pcell{4}{The students know if they chew gum. Who could still be kept blind? Which kind of blinding is that?}{1.5cm}{The person who scores the recall tests: single-blind.} \\ \hline
\end{probgrid}
\\ \hline
% ── ROW 2: completely randomized vs randomized block — THE TRAP ──────────────
\mainidea[Two designs]{Block it?} &
In a \textbf{completely randomized design}, treatments are assigned to all the units
completely \ans{at random}. In a \textbf{randomized block design}, units are first sorted into
\textbf{blocks} of units that are similar on a variable that affects the \ans{response}; then
treatments are assigned at random \ans{within} each block. Blocking takes the differences
between blocks out of the comparison of treatments.
\\
& A company tests a sports drink against water on mile times. Varsity runners are much faster
than gym-class runners, so it blocks:
\par\vspace{3pt}
\begin{center}
\begin{tikzpicture}[x=1cm, y=0.62cm]
\node[dbox] (all) at (0,0) {40\\runners};
\node[dbox] (v) at (2.2,1) {Varsity\\20};
\node[dbox] (g) at (2.2,-1) {Gym class\\20};
\node[dbox] (r1) at (4.4,1) {Random\\assignment};
\node[dbox] (r2) at (4.4,-1) {Random\\assignment};
\node[dbox] (v1) at (6.9,1.55) {10: sports drink};
\node[dbox] (v2) at (6.9,0.45) {10: water};
\node[dbox] (g1) at (6.9,-0.45) {10: sports drink};
\node[dbox] (g2) at (6.9,-1.55) {10: water};
\node[dbox] (c1) at (9.5,1) {Compare\\mile times};
\node[dbox] (c2) at (9.5,-1) {Compare\\mile times};
\draw[darr] (all) -- (v); \draw[darr] (all) -- (g);
\draw[darr] (v) -- (r1); \draw[darr] (g) -- (r2);
\draw[darr] (r1) -- (v1.west); \draw[darr] (r1) -- (v2.west);
\draw[darr] (r2) -- (g1.west); \draw[darr] (r2) -- (g2.west);
\draw[darr] (v1.east) -- (c1); \draw[darr] (v2.east) -- (c1);
\draw[darr] (g1.east) -- (c2); \draw[darr] (g2.east) -- (c2);
\end{tikzpicture}
\end{center}
\notesprompt{Read the diagram.}
\begin{probgrid}{|Y|Y|Y|}
\pcell{5}{Why block on varsity vs.\ gym class?}{1.9cm}{Experience affects mile time; comparing within a block removes that difference.} &
\pcell{6}{In each block, who decides which 10 runners get the sports drink?}{1.9cm}{Chance: a random assignment, 10 and 10.} &
\pcell{7}{A coach says, ``Block on birth month instead.'' Useful? Why?}{1.9cm}{No: birth month is not related to mile time.} \\ \hline
\end{probgrid}
\\ \hline
% ── ROW 3: matched pairs — THE CRUX ──────────────────────────────────────────
\mainidea[A block of two]{Matched pairs} &
A \textbf{matched pairs design} is a block design whose blocks have \ans{two} units.
Either pair up units that are alike on a variable related to the response, and let
\ans{chance} decide which member of each pair gets which treatment; or give each unit
\ans{both} treatments, in a random \ans{order}.
\par\vspace{3pt}
A company tests a new shoe sole that may wear more slowly. Each of 15 runners wears one shoe
with the new sole and one with the old sole for a month; a coin flip decides which foot gets
the new sole. Then the tread wear on each shoe is measured.
\notesprompt{Pairs, order, and chance.}
\begin{probgrid}{|Y|Y|Y|}
\pcell{8}{Why is this better than giving 15 runners new soles and 15 others old soles?}{2.1cm}{Runners differ in miles and weight; each runner is their own comparison.} &
\pcell{9}{Each of 12 typists types once with music and once without. How should the order be decided? Why?}{2.1cm}{Flip a coin for each typist, so practice or fatigue does not always favor one condition.} &
\pcell{10}{``Each runner wears both soles, so the pairing controls everything. Skip the coin --- put the new sole on every right foot.'' Agree? Why?}{2.1cm}{No. Right feet may wear differently; chance must still pick the foot.} \\ \hline
\end{probgrid}
\\ \hline
% ── ROW 4: guided practice — WE DO, together, fresh context ──────────────────
\mainidea[Guided practice]{The Zinc Lozenge} &
A pharmacist tests whether a zinc lozenge shortens colds. 120 adults with a new cold
volunteer; 40 are age 60 or older. Older adults tend to have longer colds. The control group
gets a look-alike lozenge with no zinc.
\par\vspace{3pt}
\begin{center}
\begin{tikzpicture}[x=1cm, y=0.62cm]
\node[dbox] (all) at (0,0) {120\\adults};
\node[dbox] (v) at (2.2,1) {60 or older\\40};
\node[dbox] (g) at (2.2,-1) {Under 60\\80};
\node[dbox] (r1) at (4.5,1) {Random\\assignment};
\node[dbox] (r2) at (4.5,-1) {Random\\assignment};
\node[dbox] (v1) at (6.9,1.55) {20: zinc};
\node[dbox] (v2) at (6.9,0.45) {20: placebo};
\node[dbox] (g1) at (6.9,-0.45) {40: zinc};
\node[dbox] (g2) at (6.9,-1.55) {40: placebo};
\node[dbox] (c1) at (9.4,1) {Compare\\days of cold};
\node[dbox] (c2) at (9.4,-1) {Compare\\days of cold};
\draw[darr] (all) -- (v); \draw[darr] (all) -- (g);
\draw[darr] (v) -- (r1); \draw[darr] (g) -- (r2);
\draw[darr] (r1) -- (v1.west); \draw[darr] (r1) -- (v2.west);
\draw[darr] (r2) -- (g1.west); \draw[darr] (r2) -- (g2.west);
\draw[darr] (v1.east) -- (c1); \draw[darr] (v2.east) -- (c1);
\draw[darr] (g1.east) -- (c2); \draw[darr] (g2.east) -- (c2);
\end{tikzpicture}
\end{center}
\notesprompt{We work these together.}
\begin{probgrid}{|Y|Y|}
\pcell{11}{Name the experimental units, the treatments, and the response variable.}{1.5cm}{120 adults with colds; zinc vs.\ placebo; days the cold lasts.} &
\pcell{12}{Why block on age?}{1.5cm}{Age affects cold length, so compare treatments within similar ages.} \\ \hline
\pcell{13}{Why give the control group a placebo instead of nothing?}{1.5cm}{To separate the zinc's effect from the placebo effect of just taking a lozenge.} &
\pcell{14}{How could this experiment be double-blind?}{1.5cm}{Neither the patients nor the staff who record cold length know who got zinc.} \\ \hline
\end{probgrid}
\\ \hline
\end{guidednotes}

% ── THE NOTES END AT GUIDED PRACTICE ─────────────────────────────────────────
% There is deliberately no "you do" here: ../homework/ is the individual practice,
% started in class in the last 8 minutes while the teacher circulates.

\end{document}
